\documentclass[10pt,a4paper]{article}
\usepackage{amsmath,amssymb,amsthm,graphicx}
\usepackage{hyperref}
\usepackage{latexsym}
\usepackage{amssymb}
\usepackage{amsthm}
\usepackage{ulem}
\usepackage{enumitem}
\setlist{itemsep=0.4pt,topsep=0pt}
\parindent 0pt
\parindent 0pt
\parskip 12pt
\newcommand{\HCF}{\mathrm{hcf}}
\newcommand{\ZZ}{\mathbb{Z}}
\newcommand{\rk}{\mathrm{rk}}
\DeclareMathOperator{\len}{length}
\DeclareMathOperator{\bl}{b}
\DeclareMathOperator{\cl}{cl}
\DeclareMathOperator{\di}{d}
\DeclareMathOperator{\Cl}{Cl}
\DeclareMathOperator{\acl}{acl}
\DeclareMathOperator{\tp}{tp}
\DeclareMathOperator{\Th}{Th}
\newcommand{\Cor}{\mathcal{C}_\mathrm{or}}
\newcommand{\Cun}{\mathcal{C}_\mathrm{un}}
\newcommand{\mons}{\mathbb{M}}
\newcommand{\strong}{\mathrm{strong}}
\newcommand{\axioms}{\Delta_{\mathrm{FSTS}}}
\newtheoremstyle{note}% hnamei
{8pt}% hSpace abovei
{14pt}% hSpace belowi
{\normalfont}% hBody fonti
{}% hIndent amounti1
{\bfseries}% hTheorem head fonti
{ }% hPunctuation after theorem headi
{.5em}% hSpace after theorem headi2
{}% hTheorem head spec (can be left empty, meaning ‘normal’)i
\theoremstyle{note}
\newtheorem{thm}{Theorem}[section]
%\newtheorem{theorem}{Theorem}
\newtheorem{lemma}[thm]{Lemma} 
\newtheorem{df}[thm]{Definition}
\newtheorem{rmk}[thm]{Remark} 
\newtheorem{prop}[thm]{Proposition}
\newtheorem{fact}[thm]{Fact}
\newtheorem{cor}[thm]{Corollary}
\newtheorem{conj}[thm]{Conjecture}
\newtheorem{ex}[thm]{Example}
\newtheorem{qu}[thm]{Question}
\newtheorem*{notation}{Notation} 
\title{A saturated unconfined system can be ordered so that it stays saturated and the small model is closed}
\author{Silvia Barbina, Enrique Casanovas}
\date{arxiv:2411.13723v3; Creative Commons Attribution 4.0 International}
\begin{document}
\maketitle
\noindent\textit{Editorial note.} The counted statement is the author's labelled environment at source0.tex lines 926-928 and its proof environment opens at line 929; the pairing is by adjacency and the retained proof range is the whole environment.\par\medskip
\noindent\textit{Editorial note.} This source is newly approved for this line; the preamble is built from the author's own package lines, macro block and theorem-environment declarations (source lines 4-65), with the author's optioned amsmath line omitted so that it cannot clash with the wrapper's amsmath.\par\medskip
\noindent\textit{Editorial note.} The source carries no label-number dump and the retained counter values are left empty, so the printed numbering of the card is the wrapper's default; the identity of the counted item is its author label.\par\medskip
\noindent\textit{Editorial note.} Metadata provenance: the source record of the pinned copy and the licence evidence file are both present; the licence is CC BY 4.0 with status supported.\par\medskip
\noindent\textit{Editorial note.} r002 note: the pre-publication gate compares the card's source\_version with the source\_version recorded in the licence evidence of the same source. This source's evidence file records exactly "2411.13723v3", while the card had carried "arXiv:2411.13723v3 (CC BY 4.0)" copied from the source record; r002 aligns the card with the evidence and changes nothing else. The evidence file itself is not modified.\par\medskip
\noindent\textit{Editorial scope.} One original result of the pinned author manuscript, counted once: existence of an ordering that keeps a saturated system saturated and closes a prescribed small submodel (author label May\_2025).\par\medskip
\noindent\textit{Allowed standard prerequisites (editorial).}\begin{itemize}
\item Saturated model. A structure of cardinality kappa is saturated when every type over a subset of smaller cardinality is realised in it.. It is the hypothesis and part of the conclusion, now for an expanded language.
\item HF-ordering with a given initial segment. Every unconfined system carries an HF-ordering having a prescribed elementary submodel as an initial segment.. It supplies the ordering the construction starts from.
\item Expansion by constants for the small model. Adding names for the elements of the small elementary submodel to the language preserves saturation of the expansion.. It is the observation that lets the author pass to a saturated model of the expanded theory.
\end{itemize}\medskip
\begin{lemma}\label{May_2025}
	Let $M$ be a saturated unconfined STS, let $A\leq M$ and assume $|A|<\kappa =|M|$. Then there is an HF-ordering $<$ of $M$ such that $M$ is saturated as an $\{R,<\}$-structure and $\cl_<(A)=A$.
\end{lemma}

\begin{proof}
	Let $<$ be an HF-ordering of $M$ having $A$ as an initial segment. Notice that the expansion $M_A= (M,a)_{a\in A}$ is saturated. Let $(M^\prime_A,<^\prime)$ be a saturated model that is elementarily equivalent to $(M_A,<)$ and has cardinality $\kappa$. For each $a\in A$ that is not a product in $M$ of two elements $a_1,a_2<a$, the sentence $\forall xy \, (x,y<a\rightarrow x\cdot y \neq a)$ holds in $(M_A,<)$. Hence, $\cl_{<^\prime}(A)=A$. Since $M_A\cong M^\prime_A$, there is some ordering $<^{\prime\prime}$ of $M$ such that $(M_A,<^{\prime\prime})\cong (M^\prime_A,<^\prime)$. Then $<^{\prime\prime}$ satisfies the requirements.
\end{proof}

\end{document}
